% ---
% id: receiver/component-catalogue/rem:bsconv
% kind: rem
% title: "Beamsplitter conventions and the port Grams"
% summary: >-
%   The symmetric beamsplitter differs from the catalogue's real one by a phase
%   on each side. The two conventions give the same port Grams whenever the
%   element is terminal with each output path inside one port and the
%   input-side phase is a scalar on the declared input space.
% prominence: standard
% status: accepted
% uses: []
% mentions: [receiver/cor:phases, receiver/thm:port-grams, def:catalogue-entries]
% assumes: [A-VAC]
% refs: []
% provenance: newly-derived
% lean:
% ---

A reader who imports a beamsplitter matrix written in another convention needs
to know which of its phase factors is free. The catalogue fixes one
convention, the core $\Bsp(x,\varphi)$ of \eqref{eq:bs}, and the interface
transcribes it (\repodoc{spec/modules/receiver/interface.yaml}).

\begin{remark}\label{rem:bsconv}
The ``symmetric'' beamsplitter convention
$\Bsp_{\mathrm{sym}}(x)\coloneqq
\left(\begin{smallmatrix}\sqrt x&i\sqrt{1-x}\\
i\sqrt{1-x}&\sqrt x\end{smallmatrix}\right)$ differs from $\Bsp(x,0)$ by
phases on \emph{both} sides, and it matters which side each one is on. Writing
$\Upsilon\coloneqq\operatorname{diag}(1,i)$ on the ordered path pair, multiplication
gives
\begin{equation}\label{eq:bssym}
  \Upsilon\,\Bsp(x,0)\,\Upsilon
  =\begin{pmatrix}1&0\\0&i\end{pmatrix}
   \begin{pmatrix}\sqrt x&\sqrt{1-x}\\ \sqrt{1-x}&-\sqrt x\end{pmatrix}
   \begin{pmatrix}1&0\\0&i\end{pmatrix}
  =\begin{pmatrix}\sqrt x& i\sqrt{1-x}\\ i\sqrt{1-x}&\sqrt x\end{pmatrix}
  =\Bsp_{\mathrm{sym}}(x) ,
\end{equation}
the sign of the $(2,2)$ entry being restored by $i\cdot i=-1$. The two factors
have different status.
\begin{itemize}[leftmargin=1.6em]
\item The \emph{left} factor is an output-side phase. If the element is
  terminal and each of its two output paths lies inside one detector port,
  Corollary~\ref{cor:phases} applies and it changes no port Gram.
\item The \emph{right} factor is an \emph{input}-side phase, which
  Corollary~\ref{cor:phases} does not cover. Consider the case in which the
  element acts directly on the inputs, each of its two input paths being
  either \tref{declared input modes} or an undeclared vacuum port
  (Assumption~\ref{ass:vac}), so that the factor acts on
  $\mathbb{C}^{\Inp}$ as a fixed diagonal unitary $\widetilde\Upsilon$,
  the vacuum columns being eliminated exactly
  (Theorem~\ref{thm:port-grams}\,(iv)). (With optics upstream of the element the
  pull-back of a diagonal phase need be neither diagonal nor unitary on
  $\mathbb{C}^{\Inp}$, and the comparison that follows is not the right one to
  make; the conclusion, to convert before declaring, still holds, and more
  strongly.) The two conventions then give
  $A_k^{\mathrm{sym}}=\widetilde\Upsilon^{\dagger}A_k\widetilde\Upsilon$ for
  every $k$: the \emph{same} family only up to that conjugation. The
  conjugation is the identity map exactly when $\widetilde\Upsilon$ is a
  scalar on $\mathbb{C}^{\Inp}$, that is, when every declared input mode
  receives the same phase from the factor, as when the element's one declared
  input path is the whole declared input space and the other input path is
  the vacuum port: a global phase on all of $\mathbb{C}^{\Inp}$ cancels in
  $A_k=\Gmat_{\Inp}^\dagger\Emd{\{k\}}\Gmat_{\Inp}$. When declared input modes
  receive different phases, as when both input paths are declared,
  $\widetilde\Upsilon$ is a genuine relative phase between them, and the two
  conventions can describe \emph{different} devices, whose Grams differ by
  that conjugation.
\end{itemize}
So the two conventions agree on the port Grams, and hence on the POVM,
whenever both conditions hold: the element is terminal with each of its
output paths inside one detector port, and the input-side factor is a scalar
on the declared input space.
Neither is automatic, which is why the catalogue fixes the one convention
$\Bsp(x,\varphi)$ of \eqref{eq:bs}, and why no catalogue entry drops a phase
it could have absorbed.
\end{remark}
